%% conference_paper.tex
%% IEEE Conference Paper — CRYSTALS-Kyber Scheduling Analysis on ESP32
%% Compressed from journal_paper.tex for 5–6 page conference format

\documentclass[conference]{IEEEtran}
\IEEEoverridecommandlockouts
\usepackage{cite}
\usepackage{amsmath,amssymb}
\usepackage{graphicx}
\usepackage{textcomp}
\usepackage{xcolor}
\usepackage{url}
\usepackage{booktabs}
\usepackage{multirow}
\usepackage{array}
\usepackage{placeins}
\usepackage{float}

\begin{document}

\title{Formal Scheduling Analysis of CRYSTALS-Kyber on Dual-Core ESP32: A DAG-Based Optimality Study
\thanks{The simulation framework, scheduling analysis code, and FIPS~203 gap analysis are available at \url{https://github.com/karans7104/PQC-SA}.}
}

\author{
\IEEEauthorblockN{Karan Sonawane}
\IEEEauthorblockA{\textit{Dept.\ of Computer Science \& Eng.} \\
\textit{IIIT Kota}, Rajasthan, India \\
2023KUCP1153}
\and
\IEEEauthorblockN{Achyut Gupta}
\IEEEauthorblockA{\textit{Dept.\ of Computer Science \& Eng.} \\
\textit{IIIT Kota}, Rajasthan, India \\
2023KUCP1158}
\and
\IEEEauthorblockN{Bhavesh}
\IEEEauthorblockA{\textit{Dept.\ of Computer Science \& Eng.} \\
\textit{IIIT Kota}, Rajasthan, India \\
2023KUCP1159}
}

\maketitle

% ====================================================================
%  ABSTRACT
% ====================================================================
\begin{abstract}
We present a formal scheduling-theoretic analysis of the dual-core CRYSTALS-Kyber implementation on the ESP32 microcontroller by Segatz and Al~Hafiz. The three KEM operations are modeled as precedence-constrained directed acyclic graphs (DAGs), and the empirical task assignment is compared against the HLFET list scheduling optimum. Under software-only timing, the schedule achieves a 0.0\% gap---formally attributed to \emph{scheduling saturation}, where the critical path consumes over 85\% of total work. A parameterized sensitivity analysis modeling ESP32 hardware accelerators (SHA/$6.1\times$, AES/$9.65\times$) reveals that acceleration breaks this optimality: the gap opens to 6.7\% for encapsulation and 2.7\% for key generation. Concurrently, a FIPS~203 compliance audit identifies nine deviations, including three critical algorithmic differences preventing interoperation with any conforming ML-KEM implementation, validated through 20/20 divergence testing across all parameter sets.
\end{abstract}

\begin{IEEEkeywords}
CRYSTALS-Kyber, ML-KEM, post-quantum cryptography, ESP32, DAG scheduling, critical path method, HLFET, FIPS~203, embedded systems.
\end{IEEEkeywords}

\IEEEpeerreviewmaketitle

% ====================================================================
%  I. INTRODUCTION
% ====================================================================
\section{Introduction}

The impending threat of quantum computers to classical public-key cryptography~\cite{shor1999} has driven NIST to standardize post-quantum algorithms. The selection of CRYSTALS-Kyber as the primary key encapsulation mechanism~\cite{kyber2021, nist2022announce} has motivated implementation efforts on resource-constrained IoT platforms. Segatz and Al~Hafiz~\cite{segatz2022} addressed this by implementing Kyber on the ESP32's dual-core Xtensa LX6 architecture, exploiting task-level parallelism with empirically derived core assignments. Prior embedded PQC work includes Kyber on ARM Cortex-M4~\cite{botros2019}, Saber on ESP32~\cite{wang2020}, and post-quantum TLS integration on ESP32~\cite{buerstinghaus2020}.

Two questions remained open. First, while the Segatz dual-core schedule was empirically effective, it was not formally analyzed against scheduling-theoretic bounds. The three KEM operations decompose naturally into precedence-constrained DAGs, enabling rigorous optimality analysis via established results~\cite{graham1969, kwok1999}. Second, the Kyber specification underwent substantial changes between its 2022 selection and its August 2024 publication as FIPS~203 (ML-KEM)~\cite{fips203}, and the 90s variant used by the implementation was excluded from standardization.

This paper makes the following contributions:
\begin{itemize}
\item A formal DAG scheduling analysis identifying \emph{scheduling saturation}---the structural condition under which the empirical schedule is provably optimal.
\item A parameterized sensitivity model predicting a 6.7\% scheduling gap for encapsulation under ESP32 hardware acceleration.
\item A systematic FIPS~203 compliance audit identifying nine gaps with severity classification, producing a reusable gap taxonomy.
\end{itemize}

% ====================================================================
%  II. BACKGROUND
% ====================================================================
\section{Background}

\subsection{CRYSTALS-Kyber}

CRYSTALS-Kyber is a lattice-based KEM operating over $R_q = \mathbb{Z}_q[X]/(X^{256}+1)$ with $q = 3329$, secured by the hardness of the Module-LWE problem~\cite{kyber2021}. It provides key generation, encapsulation, and decapsulation, with IND-CCA2 security via the Fujisaki--Okamoto transform. The 90s variant substitutes SHA-2 and AES-256-CTR for SHA-3/SHAKE. Following standardization, the algorithm was published as FIPS~203~\cite{fips203}.

\subsection{ESP32 Platform}

The ESP32 features two Xtensa LX6 cores at 240~MHz with hardware accelerators for SHA ($6.1\times$ speedup) and AES ($9.65\times$ speedup)~\cite{esp32trm, segatz2022}. Segatz and Al~Hafiz decomposed each KEM operation into sub-tasks distributed across cores using FreeRTOS semaphores.

\subsection{DAG Scheduling}

In DAG scheduling, the critical path length $L_{\mathrm{CP}}$ and total work $W$ provide lower bounds on makespan: $M^{*} \geq \max(L_{\mathrm{CP}},\, W/P)$~\cite{graham1969}. The HLFET heuristic~\cite{kwok1999, adam1974} assigns priorities by bottom-level and is optimal for two-processor unit-weight trees~\cite{coffman1972}.

% ====================================================================
%  III. METHODOLOGY
% ====================================================================
\section{Methodology}

\subsection{Simulation Framework}

We developed a host-compiled simulation framework~\cite{sonawane2026repo} by compiling the ESP32 Kyber source~\cite{pqcrystals, segatz2022repo} natively on Windows using GCC (MinGW), bypassing FreeRTOS dual-core primitives. This isolates the DAG structure from hardware-specific noise while preserving task dependencies identically to the ESP32 firmware. While absolute execution times differ from ESP32 measurements, the scheduling analysis depends on \emph{relative} weight ratios between task categories (SHA, AES, arithmetic), not absolute values. The sensitivity model (Section~III-C) explicitly parameterizes these ratios using reported ESP32 accelerator speedups, making the structural conclusions platform-transferable.

\subsection{Task Profiling and DAG Construction}

Each KEM operation was decomposed into the sub-tasks identified by Segatz: 9 for key generation, 11 for encapsulation, and 7 for decapsulation (27 total). Each was timed over 1000 iterations; the coefficient of variation was below 5\% for all tasks, confirming timing stability. Task dependency graphs were constructed from the computational flow diagrams in~\cite{segatz2022}, with each sub-task as a weighted vertex and data dependencies as directed edges. Fig.~\ref{fig:dag_combined} shows the resulting DAGs.

\begin{figure*}[!t]
\centering
\includegraphics[width=\textwidth]{dag_combined.png}
\caption{Task dependency DAGs for all three KEM operations (Kyber-512) with critical paths highlighted in red. Node colors: orange = SHA, blue = AES-CTR, green = arithmetic/NTT. Node labels show task weight in $\mu$s.}
\label{fig:dag_combined}
\end{figure*}

\subsection{Sensitivity Analysis Model}

To predict scheduling under hardware acceleration, we define accelerated weights $w_s(v) = w_{\mathrm{PC}}(v) / \alpha_s(\mathrm{type}(v))$ with per-category speedup factors: Scenario~A ($\alpha = 1$, baseline), Scenario~B (SHA/$6.1\times$, AES/$9.65\times$, from reported ESP32 measurements), and Scenario~C (SHA/$4\times$, AES/$6\times$, representing lower-tier acceleration hardware). The scheduling gap is $\mathrm{Gap}(\%) = (M_{\mathrm{Segatz}} - M_{\mathrm{HLFET}}) / M_{\mathrm{HLFET}} \times 100$.

\subsection{FIPS 203 Audit}

Each function was compared against the corresponding FIPS~203 algorithm~\cite{fips203}. Deviations were classified by severity and mapped to specific algorithm steps. Partial fixes were implemented in a separate directory, and divergence was validated through KAT testing.

% ====================================================================
%  IV. RESULTS
% ====================================================================
\section{Results}

\subsection{Task Profiling}

Table~\ref{tab:profiling} presents the compressed profiling results. In key generation, KG.1 (seed expansion) accounts for 64.9\% of total work; in encapsulation, ENC.2 (matrix generation) accounts for 48.5\%. Decapsulation contains no SHA/AES tasks and totals only 18.32~$\mu$s.

\begin{table}[!t]
\renewcommand{\arraystretch}{1.1}
\caption{Task Execution Times for Kyber-512 (Dominant Tasks)}
\label{tab:profiling}
\centering
\begin{tabular}{p{3.8cm} l r r}
\toprule
\textbf{Task} & \textbf{Op.} & \textbf{$\mu$s} & \textbf{\%} \\
\midrule
KG.1 Seed exp.\ (SHA) & KG & 249.58 & 64.9 \\
KG.2 Gen.\ matrix (AES) & KG & 76.03 & 19.8 \\
KG.3 Noise $s$ (AES) & KG & 18.30 & 4.8 \\
KG.4 Noise $e$ (AES) & KG & 18.66 & 4.8 \\
KG.5 NTT($s$) & KG & 5.66 & 1.5 \\
KG.6 NTT($e$) & KG & 5.36 & 1.4 \\
KG.7 Matmul $A\hat{s}$ & KG & 8.76 & 2.3 \\
KG.8--9 Add/pack & KG & 2.49 & 0.6 \\
\midrule
ENC.1 Unpack $pk$ & Enc & 1.06 & 0.7 \\
ENC.2 Gen.\ matrix (AES) & Enc & 74.38 & 48.5 \\
ENC.3--5 Noise (AES) & Enc & 42.38 & 27.6 \\
ENC.6 NTT($r$) & Enc & 5.65 & 3.7 \\
ENC.7--8 Matmul/inner & Enc & 12.43 & 8.1 \\
ENC.9 Inv.\ NTT & Enc & 13.14 & 8.6 \\
ENC.10--11 Add/compress & Enc & 4.40 & 2.9 \\
\midrule
DEC.1--7 (all arith.) & Dec & 18.32 & 100 \\
\bottomrule
\end{tabular}
\end{table}

\subsection{Critical Path Analysis}

Table~\ref{tab:critical_path} summarizes the critical path analysis. Key generation and decapsulation exhibit CP/Work ratios exceeding 85\%, indicating \emph{scheduling saturation}: any valid schedule placing the critical path on one core is optimal.

\begin{table}[!t]
\renewcommand{\arraystretch}{1.1}
\caption{Critical Path Analysis for Kyber-512}
\label{tab:critical_path}
\centering
\begin{tabular}{l r r r r}
\toprule
\textbf{Op.} & \textbf{Work ($\mu$s)} & \textbf{CP ($\mu$s)} & \textbf{CP/W} & \textbf{Max $S_p$} \\
\midrule
KeyGen & 384.83 & 336.86 & 87.5\% & 1.14$\times$ \\
Encaps & 153.43 & 101.29 & 66.0\% & 1.51$\times$ \\
Decaps &  18.32 &  17.75 & 96.9\% & 1.03$\times$ \\
\bottomrule
\end{tabular}
\end{table}

\subsection{Scheduling Comparison}

Table~\ref{tab:scheduling} presents the central result. Under Scenario~A, the Segatz schedule achieves 0.0\% gap for all operations---a structural consequence of scheduling saturation, not a property of the Segatz assignment in particular. For these small DAGs (7--11 nodes, 2 processors), the HLFET solution was verified to be globally optimal via exhaustive enumeration of all valid schedules. Under Scenario~B, the gap opens to 6.7\% for encapsulation (36.96 vs.\ 34.62~$\mu$s) and 2.7\% for key generation.

\begin{table*}[!t]
\renewcommand{\arraystretch}{1.1}
\caption{Scheduling Comparison: Segatz vs.\ HLFET Optimal Across Three Acceleration Scenarios (Kyber-512)}
\label{tab:scheduling}
\centering
\begin{tabular}{l l r r r r}
\toprule
\textbf{Scenario} & \textbf{Operation} & \textbf{Segatz ($\mu$s)} & \textbf{HLFET ($\mu$s)} & \textbf{Gap (\%)} & \textbf{Speedup} \\
\midrule
\multirow{3}{*}{A: PC baseline} & KeyGen & 336.86 & 336.86 & 0.0 & 1.14$\times$ \\
 & Encaps & 101.29 & 101.29 & 0.0 & 1.51$\times$ \\
 & Decaps &  17.75 &  17.75 & 0.0 & 1.03$\times$ \\
\midrule
\multirow{3}{*}{B: HW accel (SHA/$6.1\times$, AES/$9.65\times$)} & KeyGen &  61.66 &  60.05 & 2.7 & 1.25$\times$ \\
 & Encaps &  36.96 &  34.62 & 6.7 & 1.41$\times$ \\
 & Decaps &  17.75 &  17.75 & 0.0 & 1.03$\times$ \\
\midrule
\multirow{3}{*}{C: Conservative (SHA/$4\times$, AES/$6\times$)} & KeyGen &  86.32 &  86.32 & 0.0 & 1.20$\times$ \\
 & Encaps &  39.63 &  39.31 & 0.8 & 1.43$\times$ \\
 & Decaps &  17.75 &  17.75 & 0.0 & 1.03$\times$ \\
\bottomrule
\end{tabular}
\end{table*}

The mechanism is clear: hardware acceleration disproportionately shrinks SHA/AES tasks that dominate the critical path, making arithmetic tasks relatively more significant. With more balanced weights, the Segatz assignment is no longer optimal. Fig.~\ref{fig:gantt_scenariob} visualizes this gap for the encapsulation operation.

\begin{figure*}[!t]
\centering
\includegraphics[width=\textwidth]{scenario_b_gantt.png}
\caption{Gantt chart comparison under Scenario~B (ESP32 HW acceleration). Task colors: orange = SHA, blue = AES, green = arithmetic/NTT. The encapsulation panel shows the 6.7\% gap arising from suboptimal task placement in the Segatz schedule vs.\ HLFET optimal.}
\label{fig:gantt_scenariob}
\end{figure*}

\subsection{FIPS 203 Gap Analysis}

The compliance audit identified nine gaps, summarized in Table~\ref{tab:fips_gaps}. The three critical gaps all affect shared secret computation: GAP-1 omits domain separation ($G(d)$ instead of $G(d \| k)$), GAP-2 retains a removed pre-hash, and GAP-3 uses $\mathrm{KDF}(K' \| H(c))$ instead of direct $K$. Divergence testing confirmed 20/20 output divergence between the original and partially corrected versions across all three parameter sets ($k = 2, 3, 4$), with both versions independently passing round-trip consistency.

\begin{table}[H]
\renewcommand{\arraystretch}{1.1}
\caption{FIPS 203 Compliance Gaps}
\label{tab:fips_gaps}
\centering
\begin{tabular}{c l l}
\toprule
\textbf{ID} & \textbf{Description} & \textbf{Severity} \\
\midrule
1 & Missing $G(d \| k)$ domain separation & Critical \\
2 & Extra $m \leftarrow H(m)$ pre-hash & Critical \\
3 & KDF-based shared secret derivation & Critical \\
4 & No modulus check on encaps key & Moderate \\
5 & No input length/type validation & Moderate \\
6 & 90s primitives not standardized & Moderate \\
7 & Seed generated internally & Minor \\
8 & ``Kyber'' naming vs.\ ``ML-KEM'' & Minor \\
9 & Secret key layout order & Info. \\
\bottomrule
\end{tabular}
\end{table}

% ====================================================================
%  V. DISCUSSION
% ====================================================================
\section{Discussion}

The scheduling analysis reveals that the bottleneck in dual-core Kyber is structural: under software-only timing, the critical path through seed expansion and matrix generation is so dominant that scheduling decisions are trivially satisfied. The 0.0\% gap is a property of the task graph, not of the Segatz schedule. This connects to their observation that matrix generation is a candidate for optimization---our analysis provides the formal explanation.

The sensitivity model makes a falsifiable prediction: under ESP32 hardware acceleration, the HLFET-optimal assignment recovers 2.34~$\mu$s (6.7\%) for encapsulation through better interleaving of noise sampling and NTT tasks. In high-throughput IoT scenarios with thousands of encapsulations, cumulative savings become meaningful. The result demonstrates that scheduling analysis should be revisited whenever execution weights change.

The FIPS~203 findings have immediate deployment implications: any device running the current implementation cannot establish shared secrets with a conforming ML-KEM implementation. The gap taxonomy is applicable to any implementation derived from the 2022 pq-crystals codebase.

% ====================================================================
%  VI. CONCLUSION
% ====================================================================
\section{Conclusion}

This paper presented a formal DAG scheduling analysis of the dual-core Kyber ESP32 implementation, establishing that the empirical schedule is provably optimal under software timing (0.0\% gap) due to scheduling saturation, but becomes suboptimal under hardware acceleration (6.7\% for encapsulation, 2.7\% for key generation). A concurrent FIPS~203 audit identified nine compliance gaps including three critical algorithmic differences, validated through 20/20 divergence testing. Future work should validate the HLFET-optimal schedule on ESP32 hardware and migrate the symmetric primitives from the 90s variant to FIPS~203-mandated SHA-3/SHAKE.

% ====================================================================
%  ACKNOWLEDGMENT
% ====================================================================
\section*{Acknowledgment}

The authors thank Dr.\ Ritam Sharma for supervision and guidance, and IIIT Kota for providing the research environment.

% ====================================================================
%  REFERENCES
% ====================================================================
\begin{thebibliography}{16}

\bibitem{segatz2022}
F.~Segatz and M.~I.~A.~Al~Hafiz, ``Efficient implementation of CRYSTALS-KYBER key encapsulation mechanism on ESP32,'' \emph{arXiv preprint arXiv:2503.10207 [cs.CR]}, 2025.

\bibitem{kyber2021}
R.~Avanzi \emph{et~al.}, ``CRYSTALS-KYBER algorithm specifications and supporting documentation (version 3.02),'' NIST PQC Standardization, Round~3, 2021.

\bibitem{fips203}
National Institute of Standards and Technology, ``FIPS 203: Module-lattice-based key-encapsulation mechanism standard,'' Aug.\ 2024. [Online]. Available: \url{https://doi.org/10.6028/NIST.FIPS.203}

\bibitem{pqcrystals}
J.~Bos \emph{et~al.}, ``pq-crystals/kyber: Reference implementation,'' GitHub. [Online]. Available: \url{https://github.com/pq-crystals/kyber}

\bibitem{shor1999}
P.~W.~Shor, ``Polynomial-time algorithms for prime factorization and discrete logarithms on a quantum computer,'' \emph{SIAM Review}, vol.~41, no.~2, pp.~303--332, 1999.

\bibitem{graham1969}
R.~L.~Graham, ``Bounds on multiprocessing timing anomalies,'' \emph{SIAM J.\ Applied Mathematics}, vol.~17, no.~2, pp.~416--429, 1969.

\bibitem{coffman1972}
E.~G.~Coffman and R.~L.~Graham, ``Optimal scheduling for two-processor systems,'' \emph{Acta Informatica}, vol.~1, no.~3, pp.~200--213, 1972.

\bibitem{kwok1999}
Y.-K.~Kwok and I.~Ahmad, ``Static scheduling algorithms for allocating directed task graphs to multiprocessors,'' \emph{ACM Computing Surveys}, vol.~31, no.~4, pp.~406--471, 1999.

\bibitem{adam1974}
T.~L.~Adam, K.~M.~Chandy, and J.~R.~Dickson, ``A comparison of list schedules for parallel processing systems,'' \emph{Commun.\ ACM}, vol.~17, no.~12, pp.~685--690, 1974.

\bibitem{esp32trm}
Espressif Systems, ``ESP32 technical reference manual, version 4.7,'' 2022. [Online]. Available: \url{https://www.espressif.com/sites/default/files/documentation/esp32_technical_reference_manual_en.pdf}

\bibitem{botros2019}
L.~Botros, M.~J.~Kannwischer, and P.~Schwabe, ``Memory-efficient high-speed implementation of Kyber on Cortex-M4,'' in \emph{AFRICACRYPT 2019}, pp.~209--228, Springer, 2019.

\bibitem{wang2020}
B.~Wang, X.~Gu, and Y.~Yang, ``Saber on ESP32,'' in \emph{ACNS 2020}, pp.~421--440, Springer, 2020.

\bibitem{buerstinghaus2020}
K.~B\"{u}rstinghaus-Steinbach, C.~Krau\ss, R.~Niederhagen, and M.~Schneider, ``Post-quantum TLS on embedded systems: Integrating and evaluating Kyber and SPHINCS+ with Mbed TLS,'' in \emph{ASIA CCS '20}, pp.~841--852, ACM, 2020.

\bibitem{segatz2022repo}
F.~Segatz and M.~I.~A.~Al~Hafiz, ``kybesp32: Kyber on ESP32,'' GitHub. [Online]. Available: \url{https://github.com/fsegatz/kybesp32}

\bibitem{sonawane2026repo}
K.~Sonawane, A.~Gupta, and Bhavesh, ``PQC-SA: Formal scheduling analysis of CRYSTALS-Kyber on ESP32,'' GitHub, 2026. [Online]. Available: \url{https://github.com/karans7104/PQC-SA}

\bibitem{nist2022announce}
National Institute of Standards and Technology, ``NIST announces first four quantum-resistant cryptographic algorithms,'' Jul.\ 2022.

\end{thebibliography}

\end{document}
